\subsection{Fixed determinant data with degenerating ellipticity}
By \cref{thm:constant-mass-entire-determinant-profile}, for $h=\beta+\ddc u$ the normalized Gauduchon currents, their determinant measures and the sharp constant in \eqref{eq:constant-mass-determinant-profile} depend only on $\beta$. These data therefore cannot detect a loss of ellipticity of $h$. The next two theorems make this explicit on a flat torus. As in \cref{def:bmp}, $v_-$ and $v_+$ denote the lower and upper masses.
\Needspace{18\baselineskip}
\begin{theorem}\label{thm:fixed-determinant-closed-degeneration}\index{Hermitian metric!degenerating ellipticity}
Let $n\geq2$ and $X=(\mathbb C/(\mathbb Z+\mathrm i\mathbb Z))^n$, with coordinates $z_j=x_j+\mathrm iy_j$ and the flat K\"ahler form $\beta=\sum_j\mathrm dx_j\wedge\mathrm dy_j$. There
are K\"ahler forms $\omega_\varepsilon=\beta+\ddc u_\varepsilon$, $0<\varepsilon\leq1$, in the
class of $\beta$, such that the sets $\mathscr G_{\omega_\varepsilon}$ and their
families of determinant measures do not depend on $\varepsilon$, and \eqref{eq:constant-mass-determinant-profile} holds for $h=\omega_\varepsilon$ with $V=n!$ and every smooth positive probability volume $\mathrm dV$.
The minimum over $X$ of the least eigenvalue of $\omega_\varepsilon$ relative to $\beta$ equals $\varepsilon$.
All $C^k$ norms of the coefficients of $\omega_\varepsilon$ and of the
potentials $u_\varepsilon$ are bounded uniformly in $\varepsilon$, and
\[
 v_-(\omega_\varepsilon)=v_+(\omega_\varepsilon)=n!>0.
\]
The manifold has a global holomorphic volume form, so $c_1(X)=0$.
Every $\omega_\varepsilon$ is balanced and pluriclosed.
\end{theorem}
\begin{proof}
Write
$\beta_j=(\mathrm i/2)\mathrm dz_j\wedge\mathrm d\bar z_j=\mathrm dx_j\wedge\mathrm dy_j$,
so that $\beta=\sum_j\beta_j$, and put $V=\int_X\beta^n=n!$.
Set
\[
 u_\varepsilon=\frac{1-\varepsilon}{4\pi}\cos(4\pi x_1),
 \qquad a_\varepsilon=1-(1-\varepsilon)\cos(4\pi x_1).
\]
For a function of $x_1$, $\ddc u=(u''/(4\pi))\beta_1$.
Hence
\[
 \omega_\varepsilon=\beta+\ddc u_\varepsilon
 =a_\varepsilon\beta_1+\sum_{j=2}^n\beta_j.
\]
These are smooth positive closed forms with
$\varepsilon\leq a_\varepsilon\leq2-\varepsilon$ and
$\omega_\varepsilon\leq2\beta$. Their least eigenvalue relative to $\beta$
is $\min(a_\varepsilon,1)$, whose minimum over $X$ is $\varepsilon$. As $\varepsilon\to0$ they converge smoothly to a semipositive form
losing rank on $\{x_1\in\tfrac12\mathbb Z\}$.
All $C^k$ norms of the coefficients and potentials are
bounded uniformly in $\varepsilon$. \Cref{thm:constant-mass-entire-determinant-profile}, applied with this closed $\beta$ and $h=\omega_\varepsilon$, gives $\mathscr G_{\omega_\varepsilon}=\mathscr G_\beta$, the equality of the
families of determinant measures, and \eqref{eq:constant-mass-determinant-profile} with $V=n!$, for every $\varepsilon$.

Since $\omega_\varepsilon$ is closed, \cref{thm:mass-invariance-characterization} gives mass $\int_X\omega_\varepsilon^n$ for every bounded $\omega_\varepsilon$-plurisubharmonic potential. The explicit formula gives $\int_X\omega_\varepsilon^n=V$, since $a_\varepsilon$ has mean one.
Thus both extrema over bounded
potentials equal $V$. The global form
$\mathrm dz_1\wedge\cdots\wedge\mathrm dz_n$ trivializes the canonical bundle, and
closed positive forms are balanced and pluriclosed.
\end{proof}

\Needspace{24\baselineskip}
\subsection{Nonclosed pluriclosed metrics with fixed determinant data}
\begin{theorem}\label{thm:fixed-determinant-nonclosed-degeneration}\index{Pluriclosed metric!degenerating ellipticity}
On the torus $X$ of \cref{thm:fixed-determinant-closed-degeneration}, with its flat K\"ahler form $\beta$, fix $0<\tau<1/2$. There are a
nonclosed pluriclosed Hermitian metric $h_\tau$ and smooth functions $u_\varepsilon$,
$0<\varepsilon\leq1/2$, such that the forms
$\Omega_\varepsilon=h_\tau+\ddc u_\varepsilon$ are positive, and the sets $\mathscr G_{\Omega_\varepsilon}$ and their
families of determinant measures do not depend on $\varepsilon$. For every smooth positive
probability volume $\mathrm dV$ and every $p\in L^n(\mathrm dV)$,
\[
 \sup_{T\in\mathscr G_{\Omega_\varepsilon}}
       \int_X|p|\,\Gamma_0(T)
       =V_\tau^{-1/n}\|p\|_{L^n(\mathrm dV)},\qquad
 V_\tau=n!(1-\tau^2/2).
\]
The lower and upper masses of every $\Omega_\varepsilon$ equal $V_\tau$. All $C^k$ norms of the
coefficients of $\Omega_\varepsilon$ and of $u_\varepsilon$ are bounded uniformly in $\varepsilon$, while
\[
 \min_X\lambda_{\min}(\Omega_\varepsilon\mid\beta)=\varepsilon,
 \qquad \mathrm d\Omega_\varepsilon=\mathrm dh_\tau\ne0,
\]
where $\lambda_{\min}(\Omega_\varepsilon\mid\beta)$ is the least eigenvalue of $\Omega_\varepsilon$ relative to $\beta$.
Each $\Omega_\varepsilon$ is pluriclosed, and $\beta$ is a balanced metric on the same manifold,
which has $c_1(X)=0$.
\end{theorem}
\begin{proof}
With the coordinates and the forms $\beta_j$ of the proof of \cref{thm:fixed-determinant-closed-degeneration}, set $s=\sin(2\pi x_1)$ and
\[
 \xi=\frac{\mathrm i}2(\mathrm dz_1\wedge\mathrm d\bar z_2+\mathrm dz_2\wedge\mathrm d\bar z_1)
      =\mathrm dx_1\wedge\mathrm dy_2+\mathrm dx_2\wedge\mathrm dy_1,
 \qquad h_\tau=\beta+\tau s\xi.
\]
Relative to $\beta$, the matrix of $h_\tau$ in the first two coordinates is
$\left(\begin{smallmatrix}1&\tau s\\\tau s&1\end{smallmatrix}\right)$,
and its remaining diagonal entries are one, so positivity follows from $|\tau s|<1$. Moreover
\[
 \mathrm dh_\tau=2\pi\tau\cos(2\pi x_1)\,
                   \mathrm dx_1\wedge\mathrm dx_2\wedge\mathrm dy_1\ne0.
\]
Since $\xi$ has constant coefficients and $\ddc s$ is a multiple of $\mathrm dz_1\wedge\mathrm d\bar z_1$,
we have $\ddc(s\xi)=\ddc s\wedge\xi=0$. Thus the form $h'=\beta_1+\beta_2+\tau s\xi$ in the first two coordinates is pluriclosed.
Its square is a multiple of $\beta_1\wedge\beta_2$ with a coefficient depending only on $x_1$, hence closed, and its higher powers vanish. Since $h_\tau=h'+\sum_{j\geq3}\beta_j$ and the forms $\beta_j$ are closed with constant coefficients, expanding the powers of $h_\tau$ proves
\[
 \ddc(h_\tau^j)=0\quad(1\leq j\leq n-1).
\]
Moreover $h_\tau^n=n!(1-\tau^2s^2)\beta_1\wedge\cdots\wedge\beta_n$, so
\[
 \int_X h_\tau^n=n!\int_0^1(1-\tau^2\sin^2(2\pi x_1))\,\mathrm dx_1
                   =V_\tau.
\]
Use the same $u_\varepsilon$ and $a_\varepsilon$ as in
\cref{thm:fixed-determinant-closed-degeneration}, now with $\varepsilon\leq1/2$. Then
\[
 \Omega_\varepsilon=a_\varepsilon\beta_1+
       \sum_{j\geq2}\beta_j+\tau s\xi,
 \qquad a_\varepsilon=\varepsilon+2(1-\varepsilon)s^2.
\]
The first block of $\Omega_\varepsilon-\varepsilon\beta$ has
diagonal entries $2(1-\varepsilon)s^2$ and $1-\varepsilon$, and
determinant
\[
 \bigl[2(1-\varepsilon)^2-\tau^2\bigr]s^2\geq0.
\]
It is positive semidefinite, and the remaining diagonal entries of $\Omega_\varepsilon-\varepsilon\beta$ are
$1-\varepsilon$. Hence $\Omega_\varepsilon\geq\varepsilon\beta$,
and at points with $s=0$ the least eigenvalue of $\Omega_\varepsilon$ relative to $\beta$ equals $\varepsilon$.
Also $\Omega_\varepsilon\leq3\beta$, all $C^k$ norms
are bounded uniformly in $\varepsilon$, and the smooth semipositive limit as $\varepsilon\to0$ loses rank
where $s=0$. Since $\mathrm d\ddc u_\varepsilon=0$, the exterior derivative of $\Omega_\varepsilon$
is the fixed nonzero form $\mathrm dh_\tau$ computed above.

The binomial expansion in the proof of \cref{thm:constant-mass-entire-determinant-profile} applies to $\Omega_\varepsilon=h_\tau+\ddc u_\varepsilon$; it gives $\ddc(\Omega_\varepsilon^j)=0$ for $1\leq j\leq n-1$, in particular pluriclosedness. The volume of $\Omega_\varepsilon$ is $V_\tau$, either by integrating the determinant $a_\varepsilon-\tau^2s^2$ of its first block, whose mean is $1-\tau^2/2$, or by \cref{thm:mass-invariance-characterization}. That theorem also gives mass $V_\tau$ for every bounded $\Omega_\varepsilon$-plurisubharmonic potential, so both mass extrema equal $V_\tau$.
\Cref{thm:constant-mass-entire-determinant-profile}, applied with $h_\tau$ in place of $\beta$ and $h=\Omega_\varepsilon$, gives $\mathscr G_{\Omega_\varepsilon}=\mathscr G_{h_\tau}$, the equality of the
families of determinant measures, and the asserted formula for every
$p\in L^n(\mathrm dV)$. The balanced trace is fixed as well:
\[
 \int_X\Omega_\varepsilon\wedge\beta^{n-1}=n!,
\]
since the off-diagonal term has zero trace and the mean of
$a_\varepsilon$ equals one. The global holomorphic volume form gives
$c_1(X)=0$.
\end{proof}

These nonclosed metrics show that the complete determinant profile, the balanced trace and
positive mass extrema can stay fixed while the
least relative eigenvalue tends to zero. The torus is K\"ahler, however, so the examples do not settle
\cref{qu:balanced-pluriclosed-program}, which remains open.
